\documentclass[11pt]{article}
\usepackage{preamble}
\setlength{\parskip}{4pt}

\begin{document}

\daybanner{AP Statistics \textperiodcentered\ Unit 2 \textperiodcentered\ Topic 2.5 \textperiodcentered\ B: 2026-10-23 \textperiodcentered\ E: 2026-11-09 \textperiodcentered\ TEACHER KEY}{Count the Overlap Once}

\sectionbanner{CED TETHER}

{\small
\begin{itemize}[leftmargin=1.5em,itemsep=2pt,topsep=2pt]
  \item \textbf{Skill 4.B} --- Interpret statistical calculations and findings to assign meaning or assess a claim.
  \item \textbf{EU VAR-4} --- The likelihood of a random event can be quantified.
  \item \textbf{LO VAR-4.C} --- Explain why two events are (or are not) mutually exclusive. [Skill 4.B]
  \item \textbf{EK VAR-4.C.1} --- The probability that events A and B both will occur, sometimes called the joint probability, is the probability of the intersection of A and B, denoted P(A $\cap$ B).
  \item \textbf{EK VAR-4.C.2} --- Two events are mutually exclusive or disjoint if they cannot occur at the same time. So P(A $\cap$ B) = 0.
\end{itemize}
}
\textbf{Essential question:} How does an intersection determine whether two events are mutually exclusive and how their union should be counted?\par
\textbf{DOK-3 skill:} 4.B \textperiodcentered\ \textbf{FRQ pattern:} {\small\texttt{mutually-\allowbreak{}exclusive-\allowbreak{}or-\allowbreak{}not-\allowbreak{}when-\allowbreak{}the-\allowbreak{}addition-\allowbreak{}rule-\allowbreak{}double-\allowbreak{}counts}}\par
\textbf{DOK rationale:} Part (c) combines an intersection diagnosis, a union calculation, and error analysis to adjudicate a claim, explain its consequence for a reader, and construct a genuinely disjoint pair.\par

\frameworkphaseheader{Do Now (first take) $\rightarrow$ Explore (video + follow-along) $\rightarrow$ Exit (finish + turn in)}{1 $\rightarrow$ 3}{45}{%
\begin{itemize}[leftmargin=1.5em,itemsep=2pt,topsep=2pt]
  \item Board slide up at the bell. Read both students' claims without naming the overlap.
  \item Begin the combined 4.3--4.5 video follow-along at minute 5.
  \item Return to the board slide at minute 33. Circulate and have students trace which table cell each total includes.
  \item Collect at the door. Score part (c) only, E/P/I.
\end{itemize}
}{%
\begin{itemize}[leftmargin=1.5em,itemsep=2pt,topsep=2pt]
  \item Choose a claim and cite one count from the preference table.
  \item Complete the follow-along about intersections, mutually exclusive events, and conditional probability.
  \item Finish (a)--(c), revise the first take in the exit line, and turn in the sheet.
\end{itemize}
}{%
\begin{itemize}[leftmargin=1.5em,itemsep=2pt,topsep=2pt]
  \item Where do the 35 students appear inside the pizza total and inside the taco total?
  \item What does a nonzero intersection say about mutual exclusivity?
  \item What real group of students disappears if the union is reported as 1?
  \item Can you name two events here whose intersection must be empty?
\end{itemize}
}{Solo teacher. Circulate ~5 pairs. Require a table-based explanation for the subtraction and a reader consequence; do not accept `because the formula says so' alone.}

\begin{callout}{\IconWarn\ WATCH FOR}{calloutred}
\begin{itemize}[leftmargin=1.5em,itemsep=2pt,topsep=2pt]
  \item Treating two different food preferences as automatically mutually exclusive.
  \item Capping 140/120 at 1 without locating the duplicated students.
  \item Naming a proposed mutually exclusive pair that still has a nonzero table cell.
\end{itemize}
\end{callout}

\sectionbanner{SCORING PART (c) --- E / P / I}

\textbf{Expected elements}
\begin{itemize}[leftmargin=1.5em,itemsep=2pt,topsep=2pt]
  \item \textbf{rejects-mutual-exclusivity} (required) --- States the events are not mutually exclusive because the intersection contains 35 students
  \item \textbf{gives-correct-union} (required) --- Uses 105/120 or 7/8 as the probability of liking pizza or tacos
  \item \textbf{explains-double-count} (required) --- Explains that adding the two marginal totals counts all 35 students in the intersection twice
  \item \textbf{states-reader-consequence} (required) --- Says the report of 1 hides or denies the 15 students who like neither food
  \item \textbf{constructs-disjoint-pair} (required) --- Names two events in the food-preference context that cannot happen together
\end{itemize}

{\small\renewcommand{\arraystretch}{1.25}
\noindent\begin{tabular}{|p{0.5in}|p{5.9in}|}
\hline
\textbf{E} & Uses the nonzero intersection to reject mutual exclusivity, gives 105/120, explains the duplicated 35 and the false all-students message, and supplies a valid mutually exclusive pair. \\ \hline
\textbf{P} & Reaches the correct conclusion and probability but omits either the double-counting mechanism, the consequence for a reader, or a valid mutually exclusive pair. \\ \hline
\textbf{I} & Calls pizza and tacos mutually exclusive, caps 140/120 at 1, or subtracts the overlap without explaining why and gives no valid disjoint pair. \\ \hline
\end{tabular}}

\textbf{Common mistakes}
\begin{itemize}[leftmargin=1.5em,itemsep=2pt,topsep=2pt]
  \item Assuming two named food preferences must be mutually exclusive
  \item Replacing a computed value above 1 with 1 instead of diagnosing the counting error
  \item Subtracting the 35 twice, producing $70/120$
  \item Calling `likes pizza' and `does not like tacos' mutually exclusive even though 45 students satisfy both
\end{itemize}

\clearpage
\focusbanner{TODAY'S PROBLEM --- annotated}

Suppose 120 students answer two yes-or-no questions: ``Do you like pizza?'' and ``Do you like tacos?'' The table summarizes their answers. Let $P$ be the event that a randomly selected student likes pizza and $T$ the event that the student likes tacos.\par\smallskip \textbf{Jamie:} ``$P(P\text{ or }T)=80/120+60/120=140/120$. Because a probability cannot exceed 1, this must mean everyone likes at least one; I would report the answer as 1.''\par \textbf{Noor:} ``Some students may have been counted in both totals, so we need the intersection before deciding.''\par
\begin{center}
\textbf{Student food preferences (counts)}\par\smallskip
\begin{tabular}{|l|c|c|c|}
\hline
\textbf{Preference} & \textbf{Pizza} & \textbf{Not pizza} & \textbf{Total} \\ \hline
\textbf{Tacos} & 35 & 25 & 60 \\ \hline
\textbf{Not tacos} & 45 & 15 & 60 \\ \hline
\textbf{Total} & 80 & 40 & 120 \\ \hline
\end{tabular}
\end{center}
\begin{firsttakebox}{FIRST TAKE (5 min)}
Does the table support reporting everyone likes pizza or tacos? Choose a claim and cite one count.\par
\answer{Not graded. A student may initially trust the marginal totals; the finish should make the duplicated intersection visible.}
\end{firsttakebox}

\textit{Rules callout ``INTERSECTION, UNION, AND DISJOINT EVENTS (keep on the page)'' is printed on the student sheet.}\par\medskip

\dokbadge{1}~\textbf{(a)} State the number of students in $P\cap T$ and write $P(P\cap T)$ as a fraction.\par
\answer{The intersection contains 35 students, so $P(P\cap T)=35/120=7/24\approx0.292$.}

\dokbadge{2}~\textbf{(b)} Compute $P(P\cup T)$ so every student in the union is counted once. Also compute the probability that a student likes neither food.\par
\answer{Count the union once: $80+60-35=105$, so $P(P\cup T)=105/120=7/8=0.875$. The table has 15 students who like neither, so $P(\text{neither})=15/120=1/8=0.125$.}

\dokbadge{3}~\textbf{(c)} Decide whether $P$ and $T$ are mutually exclusive and which student's reasoning is warranted. Cite the 35-student intersection and your union probability, explain what Jamie's calculation would mislead a reader into believing, and name a pair of events from this context that really are mutually exclusive.\par
\answer{The events are not mutually exclusive because 35 students like both, so $P(P\cap T)=35/120>0$. Noor's reasoning is warranted. Jamie's sum counts each of those 35 students twice; subtracting that duplicated intersection gives $105/120=0.875$, not 1. Reporting 1 would mislead a reader into believing all 120 students like at least one food even though 15 like neither. One mutually exclusive pair is `likes pizza but not tacos' and `likes tacos but not pizza'; one student cannot be in both groups. Another valid pair is `likes tacos' and `does not like tacos.'}

\begin{summaryexitbox}{EXIT}
One thing the video changed about my first take: \hrulefill
\end{summaryexitbox}

\end{document}
